\section{Interpretabilidad}

La Etapa 1 ya explicó LIME y Grad-CAM; aquí se aplicaron al modelo final en vez
de repetir su teoría. Se seleccionaron cinco situaciones desde el CSV del
holdout: acierto fácil, acierto difícil, error de alta confianza, deficiencia
nutricional de potasio y una imagen de laboratorio con fondo controlado. La
selección fue programática y el archivo \texttt{cases.csv} conserva ruta,
etiqueta, predicción, confianza, entorno y fuente.

Cuando el método final fue un ensemble, se calculó Grad-CAM para la clase
predicha en la última capa convolucional de cada componente. Cada mapa se
normalizó y se combinó con los mismos pesos del soft voting. La visualización
localiza regiones que influyeron en los componentes, pero sigue siendo una
aproximación de la decisión conjunta. Grad-CAM no prueba que la región sea una
causa biológica \citep{selvaraju2017gradcam}.

\begin{figure}[H]
  \centering
  \includegraphics[height=0.82\textheight,keepaspectratio]{final_gradcam_cases.png}
  \caption{Casos del holdout y mapas Grad-CAM del método final. El color rojo
  marca mayor activación relativa dentro de cada imagen.}
  \label{fig:gradcam-final}
\end{figure}

El acierto difícil de fósforo tuvo solamente 0.283 de confianza y concentró la
activación sobre la nervadura con coloración púrpura. El modelo miró una señal
coherente y aun así dejó probabilidades cercanas para otras clases. En el caso
de potasio predicho como nitrógeno (0.269), el mapa cayó sobre el borde
amarillento y necrótico. Aquí el modelo mira el lugar correcto y aun así se
equivoca: el problema está en separar dos patrones nutricionales, no en haber
ignorado la hoja.

El error de mayor confianza cuenta otra historia. La imagen de fósforo
clasificada como potasio con 0.978 activa con fuerza el borde de la maceta y el
suelo alrededor, mientras buena parte de la hoja queda fría. Ese caso combina
fondo, encuadre vertical y amarillamiento, y muestra un atajo más preocupante
que una frontera fina entre lesiones. Los dos aciertos de roya sobre fondo
negro también producen zonas calientes cerca del contorno recortado. La red
usa puntos de roya, pero el borde artificial participa en la decisión.

LIME se ejecutó sobre el error de alta confianza cuando existió uno; de lo
contrario, sobre el acierto difícil. Se usaron 32 perturbaciones y ocho
superpíxeles visibles sobre EfficientNet-B0, el componente con mayor peso del
ensemble. A diferencia de Grad-CAM, LIME pregunta cómo cambia la salida al
ocultar segmentos de esa imagen \citep{ribeiro2016why}. Con tan pocas
perturbaciones se busca una comprobación cualitativa reproducible, no una
estimación estable para todas las hojas ni una explicación del voto completo.
La copia usada por LIME limita su lado mayor a 512 píxeles antes de generar
perturbaciones; el clasificador de todas formas la reduce a 224.

\begin{figure}[H]
  \centering
  \includegraphics[width=0.90\textwidth]{final_lime_case.png}
  \caption{Imagen final seleccionada y explicación local LIME. Los bordes
  corresponden a segmentos perturbados, no a una segmentación agronómica de la
  lesión.}
  \label{fig:lime-final}
\end{figure}

La explicación LIME se aplicó al mismo error fósforo--potasio de 0.978. Los
segmentos influyentes no quedaron encerrados en la hoja central: aparecen sobre
un fragmento de su tejido, hojas vecinas, suelo y el borde de otra maceta. Junto
con el Grad-CAM concentrado alrededor de la maceta superior, ambos métodos
indican que la salida mezcla una señal foliar débil con el contexto de captura.
La figura no se usa para corregir la etiqueta. Sirve para decidir qué tipo de
ejemplo debe añadirse: hojas con síntomas semejantes, pero fondos y encuadres
diferentes.

SHAP existe en la infraestructura histórica del repositorio, pero no se ejecutó
sobre la cabeza final. Por esa razón no se incluye una figura SHAP reciclada
como si perteneciera a esta corrida. La evidencia final de interpretabilidad
son los mapas y metadatos generados por los dos métodos anteriores.
